\documentclass[11pt]{ltjsarticle}
\usepackage[margin=25mm]{geometry}
\usepackage{booktabs}
\usepackage{amsmath}
\usepackage{xcolor}
\usepackage{pgfplots}
\pgfplotsset{compat=1.17}
\usepackage{hyperref}
\hypersetup{colorlinks=true,linkcolor=blue,urlcolor=blue}

\title{\textbf{2台のベアメタル Xinu に分散した哲学者の食事}\\
\large 哲学者5名を Pi4(AIPL JIT) と Pi3(Chandy--Misra) に 3+2 で分け、50回食事を並行実行}
\author{drone-hil / xinu-rpi4 $\times$ xinu-raz $\times$ AIPL}
\date{2026年6月5日}

\begin{document}
\maketitle

\begin{abstract}
型推論クリーンな AIPL 哲学者の食事を、\textbf{2台の実機 Xinu ノードに分散}して 50 回食事を並行実行した。
5名の哲学者を \textbf{3 + 2} に分け、\textbf{Pi4}（Cortex-A72、AIPL JIT）が 3名の環で 30 食、\textbf{Pi3}
（Cortex-A53、カーネル組込みの Chandy--Misra DiningBench）が 20 食を、\textbf{同時に}担当する。両ノードは独立に
完走し、合計 50 食を達成した。並行 wall-clock は $\max(t_{Pi4}, t_{Pi3}) = 13.5$\,s で、単一ノードで 50 食を
こなした場合（$23.4$\,s）に対し \textbf{約 $1.74\times$} の短縮を得た（Pi3 が 20 食をネイティブに $2.1$\,s で
並行吸収するため）。実装が\textbf{異種}（Pi4=AIPL 待ち行列方式・HTTP 再ポーク駆動、Pi3=組込み CM・ネイティブ）で
ある点と、これが\emph{共有フォークを持つ1つの環をメッシュ越しに調停}したものではなく\emph{2つの独立サブ問題の並行}
である点を明示する。
\end{abstract}

\section{分散の設計}
哲学者の食事は本来、隣接philがフォークを共有する\emph{結合}した環であり、N-Queens のような無損失分割はできない。
2台に「適当に分散」するにあたり、各ノードの実行系が異なる（Pi4 は任意 AIPL を JIT、Pi3 は組込み AIPL のみ）ため、
\textbf{2つの独立な食事サブ問題}を並行実行する方式を採った：
\begin{itemize}
  \item \textbf{Pi4（10.0.0.1, AIPL JIT)}：型推論クリーンな \texttt{dining3\_xinu.abcl}（単一クラス \texttt{D}・
  3名の環・順序付き獲得＋待ち行列フォーク）を \texttt{/actor/load} で投入。JIT ポンプは load ごとに quiesce する
  ため、各philが quota(=10) に達するまで \texttt{dine} を HTTP 越しに再ポークして駆動。30 食。
  \item \textbf{Pi3（192.168.3.50, 組込み CM)}：カーネル内蔵の DiningBench（mode 3 = Chandy--Misra、5名の環）を
  \texttt{/api/dining/start?mode=3\&meals=20} で起動し、\texttt{/api/dining/status} で監視。ネイティブ実行・GC 連携。
  20 食。
\end{itemize}
Mac 側のドライバが両者を\textbf{並行スレッドで同時起動}し、双方の完了を待って集計する。
分散 wall-clock $= \max(t_{Pi4}, t_{Pi3})$、総食数 $= 30 + 20 = 50$。

\section{結果}
両ノードとも正しく完走（Pi4 は各phil $10,10,10$＝完全公平、Pi3 は $n\_done=5/5$）。合計 50 食。

\begin{table}[h]\centering
\caption{2ノード分散・哲学者の食事（50 食）。}
\label{tab:dd}
\begin{tabular}{l l r r r r}
\toprule
ノード & 実装 & 食数 & 各phil & 時間 (ms) & 備考 \\
\midrule
Pi4 (A72) & AIPL JIT・3名環・待ち行列 & 30 & $10,10,10$ & 13475 & 27 pokes, GC 220ms \\
Pi3 (A53) & 組込み Chandy--Misra・5名環 & 20 & $n\_done\,5/5$ & 2092 & native, GC 連携 \\
\midrule
\multicolumn{2}{l}{\textbf{合計 / 分散 wall-clock}} & \textbf{50} & & \textbf{13475} & $=\max(\cdot)$ \\
\multicolumn{2}{l}{単一ノード 50 食（参考）} & 50 & $10{\times}5$ & 23385 & speedup $\approx 1.74\times$ \\
\bottomrule
\end{tabular}
\end{table}

\begin{figure}[h]\centering
\begin{tikzpicture}
\begin{axis}[width=11cm,height=5cm,xbar,bar width=16pt,
  xlabel={時間 (ms)}, symbolic y coords={Pi3,Pi4,Single},
  ytick=data, nodes near coords, xmin=0, xmax=27000, grid=major, grid style={gray!20}]
\addplot coordinates {(2092,Pi3) (13475,Pi4) (23385,Single)};
\end{axis}
\end{tikzpicture}\\
\footnotesize Pi3 = 20食 / Pi4 = 30食 / Single = 単一ノード50食
\caption{Pi3 は 20 食を $2.1$\,s で並行に吸収し、分散 wall-clock は Pi4 の $13.5$\,s に律速される。
単一ノード 50 食の $23.4$\,s に対し約 $1.74\times$ の短縮。}
\label{fig:dd}
\end{figure}

\section{考察}
\textbf{なぜ速くなるか。} 単一ノードでは 50 食すべてが HTTP 再ポークに律速される Pi4 上で直列に進む。分散では Pi4 の負荷が
30 食に減り、残り 20 食を Pi3 がネイティブに（再ポーク不要で）$2.1$\,s で並行処理する。よって wall-clock は
Pi4 の 30 食分（$13.5$\,s）に律速され、$23.4 \to 13.5$\,s（$\approx 1.74\times$）。
\textbf{律速の非対称性。} Pi3 は $2.1$\,s で完了後アイドルになり、$13.5-2.1 \approx 11.4$\,s 遊ぶ。Pi4 が
HTTP 再ポーク方式で重いため、より均衡した分割（Pi4 へ少なく、Pi3 へ多く）にすればさらに縮む余地がある。
\textbf{GC。} Pi4 の常駐 7 アクター（3フォーク＋3phil＋root）は GC で回収（$220$\,ms）。Pi3 の DiningBench は
GC アクターと連携してランごとに前回分を回収する（組込み）。

\section{限界（正直な但し書き）}
\begin{itemize}
  \item \textbf{異種実装}：Pi4=AIPL 待ち行列方式、Pi3=組込み Chandy--Misra。各ノードの実行系の制約
  （Pi3 は任意 AIPL を JIT できない）に由来する現実的な妥協であり、同一プログラムの対称分散ではない。
  \item \textbf{独立サブ問題}：これは\emph{共有フォークを持つ1つの環を WiFi メッシュ越しに調停}したものではなく、
  2つの独立な食事サブ問題の並行である。真のメッシュ分散（境界フォークをノード間で調停）には、受信ネットワーク
  $\to$アクターの配送フックという未実装の転送層が要る（今後）。
  \item \textbf{Pi4 は HTTP 再ポーク律速}（$\approx 460$\,ms/食）で、生の食事計算時間ではない。単一ラン。
\end{itemize}

\section{結論}
哲学者5名を 3+2 で 2 台の実機 Xinu に分け、50 回食事を\textbf{並行実行}して合計 50 食を達成、分散 wall-clock
$13.5$\,s（単一ノード $23.4$\,s に対し $\approx 1.74\times$）を実測した。Pi4 は型推論クリーンな AIPL の
待ち行列方式（完全公平）、Pi3 は組込み Chandy--Misra（$n\_done\,5/5$）で、それぞれの runtime 特性を活かした。
次段は、独立サブ問題ではなく\textbf{1つの環の境界フォークを WiFi メッシュ越しに調停}する真のメッシュ分散へ進み、
通信律速の協調コストを実測することである。

\end{document}
